\section{Background}
\label{sec:background}

\subsection{VFS Name Resolution}

The Linux VFS resolves a pathname by walking directory components, consulting
directory entries and inodes, and then dispatching ordinary file operations to
the selected filesystem~\cite{linux_vfs_docs,linux_path_lookup}. Lookup selects
the object on which later file semantics operate, while the VFS and lower
filesystem retain permission checks, reads and writes, page-cache behavior,
persistence, and consistency. The lookup/operation separation is the boundary
the paper uses for path-view policy.

\subsection{Programmable Kernel Policy}

eBPF lets the kernel run verified, bounded programs at explicit attachment
points~\cite{linux_bpf_docs,linux_bpf_verifier}. The \code{sched_ext} scheduler
class is a useful precedent because it lets BPF choose scheduling policy while
the kernel retains scheduler machinery~\cite{linux_sched_ext}. \namei applies
the same policy-versus-ownership separation to VFS name resolution: BPF supplies
a bounded path-view decision, and the kernel retains the VFS code that owns
dentries, inodes, and object lifetime. The analogy concerns the ownership split
rather than the scheduler interface.

\subsection{Namespace Construction And Filesystem Services}

Linux already offers mechanisms for alternate filesystem views. Bind mounts,
OverlayFS, projected volumes, and copied trees construct or materialize a
namespace before the application uses it~\cite{linux_mount,linux_overlayfs,kubernetes_projected_volumes}.
FUSE and stackable or custom filesystems provide a broader programmable
filesystem boundary, often by placing a service or filesystem implementation on
the operation path~\cite{linux_fuse,wrapfs,bento}. These mechanisms define the neighboring
boundaries for the path-view policy studied in the rest of the paper.
